\chapter*{Abstract}
\addcontentsline{toc}{chapter}{Abstract}
This project examines how vision-language models contribute to robot manipulation through desktop prototypes, reproduction of public policies, high-level task processing, visual completion checking, mobile manipulation in a shared scene, affordance localization, and edge deployment. It began with SO101 teleoperation and observation tools, then used action-chunking policies and physical demonstrations to examine data and action interfaces. Work subsequently focused on the public LIBERO benchmark, establishing a reproducible vision-to-action baseline with \pifive{}.

Against this baseline, the project tested Qwen for Chinese instruction decomposition, stage checking, and failure recovery. The comparisons show that translation, visual judgment, and action execution are distinct capabilities: a correct high-level goal does not guarantee that the low-level policy follows it. Later work focused on actionable locations such as drawer handles, trained low-rank adapters, and connected them to navigation, planning, and simulated joint control within one physical world. An independent Jetson environment was used to compare smaller models, quantization accuracy, cross-device outputs, and inference latency.

The current \pifive{} public-task baseline is 117/120. Frozen Chinese hierarchical control scores 28/30, compared with 29/30 for translation alone. In the latest affordance-driven pre-contact test, LoRA scores 4/4 and the base model 0/4. Handing over the full task after an additional preparatory pose succeeds in 2/4, showing that a connected interface still requires skill-stage compatibility checks. Three joint pre-contact episodes with Jetson perception and PC simulation pass, but latency in 15W mode misses the 3-second per-query target. These counts come from different tasks and conditions and cannot be combined into an overall success rate.

The report contains six technical chapters and twenty-four experimental stages, plus a chapter on engineering maintenance and delivery verification, tracing the progression from early idealized experiments through public benchmarks to physical execution in a shared simulated scene. Each chapter describes its setup, data splits, main comparisons, failures, and code locations. Appendices collect reproduction commands and media indexes. Historical and revised implementations are identified separately; in particular, the action-adapter weight-merging bug is distinguished from the capabilities of the underlying policy.

Three engineering conclusions follow. First, simple baselines such as translation alone, fixed switching, and the native policy are needed to establish whether an added module contributes. Second, two-dimensional localization, three-dimensional back-projection, planning, reaching, and complete-task success require separate acceptance checks. Third, loading a model on an edge device does not establish real-time performance. The reproducible contributions are system integration, controlled comparisons, and interface diagnosis, rather than general manipulation across unseen scenes.

This edition adds twelve sets of recorded frames, model inputs, and label examples, bringing the figure count to 72. The engineering chapter covers the shared runtime, interface validation, and 127 offline tests. The code changes are not presented as new model success rates.

\noindent\textbf{Keywords:} Vision-language-action model; Robot manipulation; Visual affordance localization; Same-scene simulation; Low-rank adaptation; Edge deployment.
